\chapter*{Заключение}
\addcontentsline{toc}{chapter}{Заключение}
\label{ch:conclusion}

В ходе выполнения лабораторной работы~\cfgLabNumber{} разработано
интерактивное веб-резюме --- одностраничное приложение, реализованное
\textbf{только базовыми технологиями} (HTML, CSS, JavaScript) без сторонних
библиотек, фреймворков и этапа сборки.

Решены все задачи методички:
\begin{itemize}
  \item страница с резюме свёрстана по типовому шаблону на семантической
        разметке HTML5 в две колонки средствами \texttt{CSS~Grid};
  \item реализовано редактирование текстовых элементов <<на месте>> через
        атрибут \texttt{contenteditable}; правки сохраняются внутри самих
        элементов и в хранилище браузера \texttt{localStorage};
  \item добавлена кнопка <<Скачать>>, экспортирующая текущую версию резюме в
        PDF средствами \texttt{window.print()} и отдельной таблицы стилей
        \texttt{@media~print} --- без единой сторонней библиотеки;
  \item все изменения элементов сопровождаются CSS-анимациями
        (\texttt{@keyframes}): подсветка при входе в правку и вспышка при
        сохранении;
  \item добавлен эффект \textbf{Material Wave} (ripple) при кликах и
        обеспечена адаптивная вёрстка для разрешений десктопа, планшета и
        смартфона;
  \item введённые данные сохраняются при перезагрузке страницы, а готовое
        приложение опубликовано на собственном домашнем сервере за обратным
        прокси \textbf{Caddy} с TLS-сертификатом Let's~Encrypt по адресу
        \url{https://resume.server34.netcraze.club}.
\end{itemize}

Дополнительные приёмы, применённые в проекте сверх явных требований:
\begin{itemize}
  \item сборка резюме из конфигурации (\texttt{config.js}) --- разделение
        данных и разметки;
  \item загрузка фотографии через \texttt{FileReader} с уменьшением размера
        средствами \texttt{<canvas>} --- чтобы изображение укладывалось в
        квоту \texttt{localStorage};
  \item конфиденциальность: в публичный репозиторий попадает только
        выдуманный персонаж, а реальные данные и фото хранятся исключительно
        в браузере пользователя;
  \item восстановление правок через \texttt{textContent} --- структурная
        защита от stored-XSS;
  \item учёт системной настройки \texttt{prefers-reduced-motion}.
\end{itemize}

Полученные навыки --- редактируемый контент, клиентское хранилище данных,
работа с файлами, генерация документов и CSS-анимации --- составляют основу
разработки современных интерактивных веб-приложений и применимы как в
проектах на чистом JavaScript, так и в работе с любыми фреймворками.

\endinput
